% !TeX spellcheck = en_US
% !TeX encoding = UTF-8
\chapter{Task 4: DDoS Filtering with eBPF/XDP}\label{ch:task4}

Detection alone does not stop an attack: once a flood is recognised, its packets must be discarded before they consume the victim's resources, and more cheaply than they are produced. We first summarise eBPF, XDP, their use cases and tools, and then present our per-source-IP rate limiter, the test bed and the measured results.

\section{eBPF and XDP}\label{sec:t4-ebpf}\label{sec:t4-xdp}
The \emph{extended Berkeley Packet Filter} (eBPF) runs small, sandboxed programs inside the Linux kernel, without kernel changes or modules \autocite{kdoc_bpf,ebpf_what}. Restricted C is compiled by clang to bytecode and JIT-compiled by the kernel. The \textbf{verifier} checks every path before loading: bounded loops, checked memory accesses (packet pointers against \texttt{data\_end}), no pointer leaks, allowed helpers only \autocite{kdoc_verifier}. State lives in \textbf{maps}, key/value stores shared with user space (hash, array, LRU, per-CPU, ring buffer, \dots) \autocite{kdoc_maps,kdoc_maphash}. The program type fixes the hook: XDP and tc for networking, kprobes and tracepoints for tracing, cgroup and LSM for security. With BTF \autocite{kdoc_btf}, libbpf relocates struct offsets at load time (CO-RE) \autocite{nakryiko_core}. Our program is small enough to need neither: it is compiled against the kernel\textquotesingle s UAPI headers and attached with iproute2.

\textbf{XDP} (eXpress Data Path) is an eBPF hook in the NIC driver, right after DMA and \emph{before} the kernel allocates an \texttt{sk\_buff} \autocite{hoiland2018xdp,ebpfdocs_xdp}. The program gets a minimal context (\texttt{struct xdp\_md}) and returns one of five actions:
\begin{itemize}
	\item \texttt{XDP\_DROP}: discard the frame; \texttt{XDP\_PASS}: hand it to the normal stack;
	\item \texttt{XDP\_TX}: send it back out of the same NIC; \texttt{XDP\_REDIRECT}: to another NIC, CPU or AF\_XDP socket;
	\item \texttt{XDP\_ABORTED}: signal an error.
\end{itemize}
Without an \texttt{sk\_buff}, drops are cheap: tens of millions of packets per second per core, the same order of magnitude as DPDK while the kernel keeps the NIC \autocite{hoiland2018xdp}. Cloudflare measured XDP drops faster than iptables \texttt{raw} and tc \autocite{cf_drop10m}. iptables/nftables stay the tool for stateful firewalling; XDP suits early decisions on huge volumes, i.e.\ DDoS. Three attach modes exist \autocite{ebpfdocs_xdp,cilium_bpfref}:
\begin{itemize}
	\item \textbf{native}: in the driver's NAPI loop, the fast path (\texttt{veth}, \texttt{virtio\_net}, most fast NICs);
	\item \textbf{generic}: after \texttt{sk\_buff} allocation, driver-independent and slow, good for testing;
	\item \textbf{offloaded}: on a SmartNIC.
\end{itemize}

\section{Use Cases and Tools}\label{sec:t4-usecases}\label{sec:t4-tools}
Main use cases of XDP:
\begin{itemize}
	\item \textbf{DDoS mitigation}: Cloudflare's \emph{Gatebot} detects attacks, \emph{L4Drop} compiles rules into XDP, \emph{flowtrackd} tracks TCP state \autocite{cf_gatebot,cf_l4drop,cf_flowtrackd}.
	\item \textbf{Load balancing}: Katran (IP-in-IP) and Unimog (GUE, UDP-based) forward packets with \texttt{XDP\_TX} after encapsulation \autocite{meta_katran,katran_github,cf_unimog}.
	\item \textbf{Firewalling}: Cilium uses XDP for CIDR pre-filtering and kube-proxy replacement \autocite{cilium_kpr,cilium_bpfref}.
	\item \textbf{Monitoring}: counting, sampling and capture at line rate (\texttt{xdpdump}).
	\item \textbf{IDS/IPS}: Suricata bypasses flows it has already inspected \autocite{suricata_ebpf}.
	\item \textbf{AF\_XDP}: zero-copy delivery of selected packets to user space \autocite{kdoc_afxdp}.
\end{itemize}
\Cref{tab:t4-tools} lists the tools relevant here. We used iproute2 to attach our program and bpftool to inspect it; Task~5 covers xdp-tools.

\begin{table}[htbp]
	\centering\footnotesize
	\caption{Survey of eBPF/XDP tools.}\label{tab:t4-tools}
	\begin{tabularx}{\linewidth}{>{\raggedright\arraybackslash}p{4.3cm}>{\raggedright\arraybackslash}p{2.2cm}X}
		\toprule
		Tool & Type & Description \\
		\midrule
		xdp-tools \autocite{xdptools_github} & XDP utilities & libxdp, \texttt{xdp-loader}, \texttt{xdp-filter}, \texttt{xdpdump} (\cref{ch:task5}). \\
		xdp-tutorial \autocite{xdp_tutorial} & Teaching & Lessons; basis of our program structure. \\
		libbpf \autocite{libbpf_github}, bpftool \autocite{kdoc_bpf} & Library, CLI & libbpf: BPF loading library used by \texttt{ip} and bpftool; bpftool: inspects programs and maps. \\
		iproute2 & CLI & Attaches our XDP object (\texttt{ip link set \dots\ xdpgeneric obj}). \\
		BCC \autocite{bcc_github}, bpftrace \autocite{bpftrace_github} & Tracing & Python tools; one-liner tracing. \\
		Cilium \autocite{cilium_kpr}, Katran \autocite{katran_github} & Networking & Kubernetes networking; XDP load balancer. \\
		\bottomrule
	\end{tabularx}
\end{table}

\section{Our Filter: a Per-Source-IP Rate Limiter}\label{sec:t4-design}
A volumetric flood such as a SYN flood sends far more packets per second from one source than any legitimate client. Our XDP program therefore counts packets per IPv4 source address in a fixed window of one second and drops every packet above a limit of 100 packets per second. We kept the program deliberately small (78 lines): the limit is a compile-time constant, and all state lives in a single BPF map.

\paragraph{Decision path.} For every frame, the program (\cref{lst:t4-prog}) first checks that the Ethernet and IPv4 headers lie within \texttt{data\_end}; without this bounds check the verifier rejects the program. Frames that are not IPv4 (ARP, IPv6) are passed. For IPv4, the program looks up the source address in \texttt{rl\_state} and creates an entry on the first packet. If the entry's window has expired, it is reset. The packet count is then incremented atomically, and a packet whose count exceeds the limit gets \texttt{XDP\_DROP}; all others get \texttt{XDP\_PASS}. Each entry also counts passed and dropped packets, so the map shows the effect per source.

\lstinputlisting[style=bpfc,float=tbp,firstnumber=35,firstline=35,lastline=76,caption={The XDP program in \texttt{xdp\_ratelimit.bpf.c} (includes and map definition above it omitted).},label=lst:t4-prog]{../xdp/ratelimit/xdp_ratelimit.bpf.c}

\paragraph{Map choice.} \texttt{rl\_state} is an \texttt{LRU\_HASH} with 65,536 entries (key: the four bytes of the source address; value: window start, count, passed, dropped). At capacity the least recently used source is evicted, whereas a plain hash map would refuse new entries, which an attacker could exploit. It is \emph{shared}, not per-CPU, on purpose: with receive-side scaling, one source's packets can be spread over several CPUs, and per-CPU counters would raise the effective limit to $\mathrm{limit} \times \#\mathrm{CPUs}$. The price is an atomic \texttt{fetch\_add} (which requires the BPF v3 instruction set, \texttt{-mcpu=v3}) and a benign race on the window reset, in which at worst a few extra packets pass.

\paragraph{Build and loading.} \texttt{make} compiles the object with \texttt{clang -O2 -g -target bpf -mcpu=v3} against the kernel's UAPI headers and libbpf's \texttt{bpf\_helpers.h}. No user-space loader is needed: iproute2 attaches the object (\texttt{ip link set dev veth0 xdpgeneric obj xdp\_ratelimit.bpf.o sec xdp}), and bpftool dumps the map; thanks to the BTF debug information (\texttt{-g}), it prints the map entries with field names. The verifier accepted the program: \texttt{ip} attached it without error.

\section{Test Bed}\label{sec:t4-demo}
To test the filter without touching a physical interface, we use a single network namespace \texttt{atk}, connected to the host by a veth pair. The namespace has two addresses on its interface: 10.0.0.2 acts as the attacker and 10.0.0.3 as a benign client. The host (10.0.0.1, interface \texttt{veth0}) is the victim, and the rate limiter is attached to \texttt{veth0} in generic (skb) mode. The attacker runs a SYN flood to port~80 (\texttt{hping3 -S -p 80 --flood}) for 5\,s. We count the packets that reach the host's IP layer, i.e.\ that passed XDP, with the kernel counter \texttt{IpInReceives} (\texttt{nstat}, reset before each round), and read the exact per-source counts from the map. The commands are typed by hand (they are listed in \texttt{xdp/ratelimit/README.md}): a baseline flood without XDP, attaching the program with \texttt{ip}, the same flood while the benign client pings the victim five times per second, and a dump of \texttt{rl\_state} with bpftool. The limit is 100 packets per second and source.

\section{Results}\label{sec:t4-results}
\Cref{tab:t4-results} summarises both rounds. Without XDP, the attacker sent 372,537 packets in 5\,s, and the host's IP layer received 372,564 packets, i.e.\ the whole flood plus 27 packets of other host traffic (\cref{fig:ss-task4_baseline.png}). With the rate limiter, the attacker sent 697,064 packets, and the IP layer received only 525 (\cref{fig:ss-task4_flood_xdp.png}): the 500 flood packets that the limiter let through (100 per second for 5\,s, as configured) and the 25 benign pings, all of which were answered (0\,\% loss). The attacker could send almost twice as many packets with the filter active; we attribute this to the CPU time the host spent on every SYN in the baseline, since attacker and victim share the same machine.

\begin{table}[htbp]
	\centering\small
	\caption{Results of the two flood rounds (limit 100 packets/s per source, 5\,s flood; values from \texttt{task4\_terminal\_output.txt}).}\label{tab:t4-results}
	\begin{tabular}{lrrr}
		\toprule
		Round & Packets sent & Packets at victim (IP layer) & Benign ping loss \\
		\midrule
		Baseline, no XDP & 372,537 & 372,564 & -- \\
		XDP, generic mode (\texttt{ip}) & 697,064 & 525 & 0\,\% (25/25) \\
		\bottomrule
	\end{tabular}
\end{table}

\paragraph{Map.} The dump of \texttt{rl\_state} after the attack (\cref{fig:ss-task4_map_dump.png}) shows one entry per source. The attacker \texttt{[10,0,0,2]} has 500 packets passed and 696,564 dropped, which together are exactly the 697,064 packets hping3 sent. The benign client \texttt{[10,0,0,3]} has 25 passed and none dropped. The limit held exactly (500 in 5\,s). Because the entry is shared rather than per-CPU, this also holds when one source's packets are processed on several CPUs.

\screenshot[0.8\linewidth]{task4_baseline.png}{Baseline without XDP: 372,537 SYN packets sent, 372,564 packets reach the host's IP layer.}
\screenshot[0.8\linewidth]{task4_flood_xdp.png}{Flood with the rate limiter attached: 697,064 packets sent, 525 reach the host's IP layer (500 flood packets and 25 pings); the benign ping has 0\,\% loss.}
\screenshot[0.7\linewidth]{task4_map_dump.png}{\texttt{bpftool map dump name rl\_state}: 500 passed and 696,564 dropped for the attacker, 25 passed and 0 dropped for the benign client.}

\section{Problems and Fixes}\label{sec:t4-problems}
\emph{The firewall hid the flood.} In our first run we counted the TCP segments that reached the host (\texttt{TcpInSegs}). The counter stayed near zero in both rounds, also without XDP, although hundreds of thousands of SYNs were sent. The reason is the host firewall (ufw, default policy: deny incoming): netfilter drops the SYNs before TCP sees them, so \texttt{TcpInSegs} only counted the host's own traffic. We therefore count with \texttt{IpInReceives}, which the kernel increments when a packet enters the IP layer, before netfilter, and which thus shows exactly what passed XDP.

\section{Limitations and Bridge to the ML Model}\label{sec:t4-limits}\label{sec:t4-bridge}
(1)~\textbf{Spoofed sources.} Per-source limiting assumes that the source IP identifies the sender. A SYN flood with random spoofed sources sends one or two packets per address and is never limited; it also churns the LRU map. Real systems add aggregate limits, SYN cookies and signature rules as in L4Drop \autocite{cf_l4drop}. (2)~\textbf{Window boundary.} A fixed window lets up to $2\times$limit packets pass around a boundary; a token bucket would smooth this. (3)~\textbf{Shared addresses.} Users behind one carrier-grade NAT share one budget. (4)~\textbf{IPv4 only.} IPv6 is passed unfiltered. (5)~\textbf{Ingress only, small test bed.} XDP sees only received packets, and our tests used veth on one host with a single attacker. (6)~\textbf{Fixed limit.} The limit is a compile-time constant; changing it means recompiling and re-attaching the program.

The program has no interface for a detector yet; this is the natural next step. A second hash map with blocked source addresses, checked before the rate limit, would let a user-space classifier (our Lab~3 model or a Lab~4 detector) act on its decision: it inserts a flagged source with \texttt{bpf\_map\_update\_elem()} or \texttt{bpftool map update}, and from the next packet on that source is dropped in the driver. This is the division of labour Cloudflare uses at scale: Gatebot detects in user space, L4Drop enforces in XDP \autocite{cf_gatebot,cf_l4drop}. The classifier's decision latency (\cref{sec:cost}) then matters once per flow, not per packet.
